\section{What has closed and what remains}
\label{sec:tpc41-next-gate}

The endpoint can now be stated with fewer moving parts.

\subsection{Unconditional advances}

The following statements are proved without a new conjectural input:
\begin{enumerate}[label=(\roman*)]
 \item the complete signed same--row unequal--alias sector in the minimal
 bank is bounded by the existing atomic diagonal;
 \item the full alias DFT deletes that sector exactly;
 \item the physical terminal bank is reduced to the cross--row form
 \(\cF_L+\cF_R\);
 \item every affine target pair in that form is algebraically
 nonproportional;
 \item the form admits an exact Hilbert--valued one--M\"obius
 compression; and
 \item almost all admissible triple--prime CRT moduli give a raw scalar
 progression energy smaller than the coherent scale by
 \((\log X)^\theta\), including the smooth terminal weight.
\end{enumerate}

The sixth item is a useful quantitative specialization with a stronger
exceptional--density ledger than the elementary classical benchmark.  It
should not be read as the first known scalar logarithmic saving in this
range.  The first five items remove false alternatives and specify the
coefficient class that a stronger theorem must address.

\subsection{The next analytic theorem}

The exact continuation for the selected minimal folded route is the
one--sided estimate
\begin{equation}
 \operatorname{Re}(\cF_L+\cF_R)
 \ll X^{o(1)}Q^2JK_B.
 \label{eq:tpc41-next-weakest}
\end{equation}
Alternatively, the original equality column may be attacked directly
through the distinct sufficient vector inequality
\eqref{eq:tpc41-Hilbert-equality-target}, without first bounding the
entire folded energy.  A stronger, cleaner target is the full trace inequality
\eqref{eq:tpc41-Hilbert-trace-target}.

There are three concrete routes, with distinct failure modes:
\begin{enumerate}[label=(\alph*)]
 \item \textbf{Vectorize the progression theorem.}  Prove a Bessel or
 large--sieve inequality for the actual TPC--36 projective coefficient
 class.  A scalar residue variance is insufficient unless the upgrade is
 uniform under arbitrary \(\ell^2\) dualization of the output coordinates.
 \item \textbf{Control the localized convolution defect.}  Compare
 \(\Lambda_{L_0}*\mu_{D_0}*\mathbf1_{\cJ}\), with its opposite--row masks,
 to a localized \(\Lambda\) model in a Hilbert norm.  The full identity
 \(\Lambda*\mu*\mathbf1=\Lambda\) identifies the hoped--for main term but
 supplies no boundary estimate.
 \item \textbf{Average the selectable CRT triple.}  The auxiliary primes
 are chosen inside admissible intervals.  A nonnegative average of the
 \emph{complete physical vector energy} over those triples would select
 at least one good triple.  The scalar theorem of
 \cref{thm:tpc41-raw-triple-prime} shows that this family is arithmetically
 useful, but an average of scalar residue energies is not yet an average
 of \(\cF_L+\cF_R\).
\end{enumerate}

The routes may be combined: a triple average could first yield a vector
bound for most projective layers, while a convolution estimate treats the
exceptional masks.  Conversely, a rigorous obstruction showing that one
route cannot reach \eqref{eq:tpc41-next-weakest} would itself delimit the
method class and justify returning to another branch.

\subsection{Claim boundary}

This paper does \emph{not} prove:
\begin{enumerate}[label=(\roman*)]
 \item cancellation for an individual long shift \(rM\);
 \item a diagonal--scale scalar progression variance;
 \item the Hilbert--valued physical estimate
 \eqref{eq:tpc41-Hilbert-equality-target};
 \item cancellation in the cross--row four--M\"obius form;
 \item a breach of sieve parity;
 \item a Hardy--Littlewood prime--pair asymptotic; or
 \item infinitely many twin primes.
\end{enumerate}
The result is a forward reduction plus a strict scalar logarithmic saving.
The remaining distance is explicitly the factor
\(K/(\log X)^\theta\) in the scalar shadow and the scalar--to--physical
Hilbert upgrade in the true row gate.
